\documentclass[12pt]{article}

\usepackage[spanish,es-noshorthands]{babel}
\usepackage[utf8]{inputenc}
\usepackage[T1]{fontenc}
\usepackage{amsmath,amssymb,mathtools}
\usepackage{geometry}
\usepackage{tikz}
\usetikzlibrary{arrows.meta,positioning,bending}

\geometry{margin=2.5cm}

\title{Ejercicios de Cadenas de Markov}
\author{}
\date{}

\begin{document}

\maketitle

\section*{Ejercicio 1}

Considere la cadena de Markov con espacio de estados
\[
S=\{1,2,3\},
\]
y matriz de transición
\[
P=
\begin{pmatrix}
\frac12 & \frac14 & \frac14\\[4pt]
\frac13 & 0 & \frac23\\[4pt]
\frac12 & \frac12 & 0
\end{pmatrix}.
\]

\subsection*{Diagrama de transición}

De la matriz se observa que:

\[
p_{11}=\frac12,\quad p_{12}=\frac14,\quad p_{13}=\frac14,
\]
\[
p_{21}=\frac13,\quad p_{22}=0,\quad p_{23}=\frac23,
\]
\[
p_{31}=\frac12,\quad p_{32}=\frac12,\quad p_{33}=0.
\]

Entonces el diagrama de transición es:

\begin{center}
\begin{tikzpicture}[>=Latex, node distance=3cm, auto]
\node[circle,draw] (1) {$1$};
\node[circle,draw,right of=1] (2) {$2$};
\node[circle,draw,below=2cm of 1] (3) {$3$};

\draw[->] (1) edge[loop above] node {$\frac12$} (1);
\draw[->] (1) edge[bend left] node {$\frac14$} (2);
\draw[->] (1) edge[bend right] node[left] {$\frac14$} (3);

\draw[->] (2) edge[bend left] node {$\frac13$} (1);
\draw[->] (2) edge[bend left] node[right] {$\frac23$} (3);

\draw[->] (3) edge[bend right] node[left] {$\frac12$} (1);
\draw[->] (3) edge[bend left] node {$\frac12$} (2);
\end{tikzpicture}
\end{center}

\subsection*{Cálculo de \(P(X_1=3,X_2=2,X_3=1)\)}

Nos dicen que
\[
P(X_1=1)=P(X_1=2)=\frac14.
\]
Entonces
\[
P(X_1=3)=1-\frac14-\frac14=\frac12.
\]

Ahora usamos la propiedad de Markov:
\[
P(X_1=3,X_2=2,X_3=1)
=
P(X_1=3)\,P(X_2=2\mid X_1=3)\,P(X_3=1\mid X_2=2).
\]

De la matriz de transición se tiene:
\[
P(X_2=2\mid X_1=3)=p_{32}=\frac12,
\qquad
P(X_3=1\mid X_2=2)=p_{21}=\frac13.
\]

Sustituyendo:
\[
P(X_1=3,X_2=2,X_3=1)
=
\frac12\cdot\frac12\cdot\frac13
=
\frac{1}{12}.
\]

Por lo tanto,
\[
\boxed{P(X_1=3,X_2=2,X_3=1)=\frac{1}{12}}.
\]

\newpage
\section*{Ejercicio 2}

Se quiere calcular la probabilidad de que la cadena se absorba en la clase recurrente
\[
R_1=\{1,2\},
\]
suponiendo que \(X_0=3\).

Sea
\[
h_i=P_i(\text{absorberse en }R_1).
\]

Como los estados \(1\) y \(2\) ya están dentro de \(R_1\), entonces
\[
h_1=h_2=1.
\]

Y como los estados \(5,6,7\) pertenecen a otra clase recurrente, se tiene
\[
h_5=h_6=h_7=0.
\]

Ahora planteamos las ecuaciones para los estados transitorios.

Desde el estado \(3\), la cadena va a \(1\) con probabilidad \(\frac12\) y a \(4\) con probabilidad \(\frac12\), por lo que:
\[
h_3=\frac12(1)+\frac12 h_4.
\]

Desde el estado \(4\), la cadena va a \(1\) con probabilidad \(\frac14\), a \(3\) con probabilidad \(\frac14\) y a \(7\) con probabilidad \(\frac12\). Entonces:
\[
h_4=\frac14(1)+\frac14 h_3+\frac12(0).
\]

Así, el sistema queda:
\[
h_3=\frac12+\frac12 h_4,
\qquad
h_4=\frac14+\frac14 h_3.
\]

Sustituyendo la segunda ecuación en la primera:
\[
h_3=\frac12+\frac12\left(\frac14+\frac14 h_3\right).
\]

Desarrollando:
\[
h_3=\frac12+\frac18+\frac18 h_3
=\frac58+\frac18 h_3.
\]

Pasamos el término con \(h_3\) al lado izquierdo:
\[
h_3-\frac18 h_3=\frac58.
\]

\[
\frac78 h_3=\frac58.
\]

Por lo tanto,
\[
h_3=\frac{5}{7}.
\]

Entonces, la probabilidad de absorberse en \(R_1\), empezando en \(3\), es
\[
\boxed{\frac{5}{7}}.
\]

\newpage
\section*{Ejercicio 3}

Ahora se pide calcular el tiempo esperado de absorción, empezando en \(X_0=3\).

Sea
\[
T=\min\{n\geq 0: X_n\in R_1\cup R_2\},
\]
es decir, el tiempo en que la cadena entra por primera vez a alguna clase recurrente.

Definimos
\[
t_i=E_i[T].
\]

Como en los estados recurrentes la absorción ya ocurrió, se tiene:
\[
t_1=t_2=t_5=t_6=t_7=0.
\]

Solo faltan los estados transitorios \(3\) y \(4\).

Desde el estado \(3\):
\[
t_3=1+\frac12 t_1+\frac12 t_4.
\]
Como \(t_1=0\), queda
\[
t_3=1+\frac12 t_4.
\]

Desde el estado \(4\):
\[
t_4=1+\frac14 t_1+\frac14 t_3+\frac12 t_7.
\]
Como \(t_1=t_7=0\), queda
\[
t_4=1+\frac14 t_3.
\]

Entonces el sistema es
\[
t_3=1+\frac12 t_4,
\qquad
t_4=1+\frac14 t_3.
\]

Sustituyendo la segunda ecuación en la primera:
\[
t_3=1+\frac12\left(1+\frac14 t_3\right).
\]

Desarrollando:
\[
t_3=1+\frac12+\frac18 t_3
=\frac32+\frac18 t_3.
\]

Pasando el término con \(t_3\) al lado izquierdo:
\[
t_3-\frac18 t_3=\frac32.
\]

\[
\frac78 t_3=\frac32.
\]

Por lo tanto,
\[
t_3=\frac32\cdot\frac87=\frac{12}{7}.
\]

Así, el tiempo esperado de absorción es
\[
\boxed{E[T\mid X_0=3]=\frac{12}{7}}.
\]

\newpage
\section*{Ejercicio 4}

Se pide calcular el tiempo medio de recurrencia al estado \(1\), suponiendo que \(X_0=1\).

Sea
\[
R=\min\{n\geq 1:X_n=1\}.
\]

La matriz de transición es
\[
P=
\begin{pmatrix}
\frac14 & \frac12 & \frac14\\[4pt]
\frac13 & 0 & \frac23\\[4pt]
\frac12 & 0 & \frac12
\end{pmatrix}.
\]

Primero calculamos el tiempo esperado para llegar a \(1\) empezando desde \(2\) y desde \(3\).

Sea
\[
m_2=E_2[\tau_1],\qquad m_3=E_3[\tau_1],
\]
donde
\[
\tau_1=\min\{n\geq 0:X_n=1\}.
\]

Desde el estado \(3\):
\[
m_3=1+\frac12\cdot 0+\frac12 m_3.
\]

Entonces:
\[
m_3-\frac12 m_3=1
\quad\Longrightarrow\quad
\frac12 m_3=1
\quad\Longrightarrow\quad
m_3=2.
\]

Desde el estado \(2\):
\[
m_2=1+\frac13\cdot 0+\frac23 m_3.
\]

Sustituyendo \(m_3=2\):
\[
m_2=1+\frac23(2)=1+\frac43=\frac73.
\]

Ahora calculamos el retorno al estado \(1\).

Partiendo de \(1\):

\begin{itemize}
    \item con probabilidad \(\frac14\), regresa inmediatamente en 1 paso;
    \item con probabilidad \(\frac12\), va a \(2\) y luego tarda en promedio \(m_2\) pasos más;
    \item con probabilidad \(\frac14\), va a \(3\) y luego tarda en promedio \(m_3\) pasos más.
\end{itemize}

Por eso,
\[
E[R\mid X_0=1]
=
\frac14(1)+\frac12(1+m_2)+\frac14(1+m_3).
\]

Sustituyendo:
\[
E[R\mid X_0=1]
=
\frac14+\frac12\left(1+\frac73\right)+\frac14(1+2).
\]

\[
E[R\mid X_0=1]
=
\frac14+\frac12\left(\frac{10}{3}\right)+\frac34.
\]

\[
E[R\mid X_0=1]
=
\frac14+\frac{5}{3}+\frac34
=
1+\frac53
=
\frac83.
\]

Por lo tanto,
\[
\boxed{E[R\mid X_0=1]=\frac83}.
\]

\newpage
\section*{Ejercicio 5}

Considere la cadena de dos estados con matriz de transición
\[
P=
\begin{pmatrix}
1-a & a\\
b & 1-b
\end{pmatrix}.
\]

\subsection*{Distribución estacionaria}

Sea
\[
\pi=(\pi_0,\pi_1).
\]
Buscamos que
\[
\pi=\pi P,
\qquad
\pi_0+\pi_1=1.
\]

De la ecuación \(\pi=\pi P\) se obtiene:
\[
\pi_0=\pi_0(1-a)+\pi_1 b.
\]

Reacomodando:
\[
\pi_0 a=\pi_1 b.
\]

Junto con
\[
\pi_0+\pi_1=1,
\]
se obtiene:
\[
\pi_0=\frac{b}{a+b},
\qquad
\pi_1=\frac{a}{a+b}.
\]

Entonces la distribución estacionaria es
\[
\boxed{\pi=\left(\frac{b}{a+b},\frac{a}{a+b}\right)}.
\]

\subsection*{Demostrar que \(r_0=\frac{1}{\pi_0}\)}

Sea \(r_0\) el tiempo medio de recurrencia del estado \(0\).

Primero calculamos el tiempo esperado para regresar a \(0\) empezando desde \(1\). Sea
\[
h_1=E_1[\tau_0].
\]

Desde \(1\), con probabilidad \(b\) se va a \(0\), y con probabilidad \(1-b\) permanece en \(1\). Entonces:
\[
h_1=1+b(0)+(1-b)h_1.
\]

De aquí:
\[
h_1-(1-b)h_1=1
\quad\Longrightarrow\quad
bh_1=1
\quad\Longrightarrow\quad
h_1=\frac1b.
\]

Ahora empezando desde \(0\):

\begin{itemize}
    \item con probabilidad \(1-a\), se queda en \(0\), así que retorna en 1 paso;
    \item con probabilidad \(a\), va a \(1\), usa un paso y luego tarda en promedio \(h_1\) pasos más en volver.
\end{itemize}

Entonces:
\[
r_0=(1-a)(1)+a(1+h_1).
\]

Sustituyendo \(h_1=\frac1b\):
\[
r_0=(1-a)+a\left(1+\frac1b\right).
\]

\[
r_0=1+\frac{a}{b}
=\frac{a+b}{b}.
\]

Pero como
\[
\pi_0=\frac{b}{a+b},
\]
entonces
\[
\frac1{\pi_0}=\frac{a+b}{b}.
\]

Por lo tanto,
\[
\boxed{r_0=\frac1{\pi_0}}.
\]

\subsection*{¿Cuándo la distribución estacionaria también es límite?}

Si la cadena es irreducible y aperiódica, entonces la distribución estacionaria también es la distribución límite.

En este caso, cuando \(a>0\) y \(b>0\), la cadena es irreducible. Además, si no ocurre el caso periódico extremo \(a=b=1\), entonces la cadena es aperiódica.

Por lo tanto, en los casos del ejercicio, la distribución estacionaria también es distribución límite.

\[
\boxed{\text{La distribución estacionaria también es distribución límite.}}
\]

\newpage
\section*{Ejercicio 6}

Considere la cadena con matriz
\[
P=
\begin{pmatrix}
\frac14 & \frac12 & \frac14\\[4pt]
\frac13 & 0 & \frac23\\[4pt]
\frac12 & 0 & \frac12
\end{pmatrix}.
\]

\subsection*{¿Es irreducible?}

Sí es irreducible, porque todos los estados se comunican entre sí.

Por ejemplo:

\begin{itemize}
    \item desde \(1\) se puede ir a \(2\) y a \(3\);
    \item desde \(2\) se puede ir a \(1\) y a \(3\);
    \item desde \(3\) se puede ir a \(1\), y de \(1\) se puede llegar a \(2\).
\end{itemize}

Entonces,
\[
\boxed{\text{La cadena es irreducible.}}
\]

\subsection*{¿Es aperiódica?}

Sí, porque
\[
p_{11}=\frac14>0.
\]

Eso significa que desde el estado \(1\) se puede regresar a \(1\) en un paso, así que su periodo es \(1\). Como la cadena es irreducible, todos los estados tienen el mismo periodo.

Por lo tanto,
\[
\boxed{\text{La cadena es aperiódica.}}
\]

\subsection*{Distribución estacionaria}

Sea
\[
\pi=(\pi_1,\pi_2,\pi_3).
\]

Queremos resolver
\[
\pi=\pi P,
\qquad
\pi_1+\pi_2+\pi_3=1.
\]

De \(\pi=\pi P\) se obtiene:
\[
\pi_2=\frac12\pi_1.
\]

También:
\[
\pi_3=\frac14\pi_1+\frac23\pi_2+\frac12\pi_3.
\]

Pasamos \(\frac12\pi_3\) al lado izquierdo:
\[
\frac12\pi_3=\frac14\pi_1+\frac23\pi_2.
\]

Sustituyendo \(\pi_2=\frac12\pi_1\):
\[
\frac12\pi_3=\frac14\pi_1+\frac23\left(\frac12\pi_1\right)
=\frac14\pi_1+\frac13\pi_1
=\frac{7}{12}\pi_1.
\]

Entonces:
\[
\pi_3=\frac76\pi_1.
\]

Ahora usamos que
\[
\pi_1+\pi_2+\pi_3=1.
\]

Sustituyendo:
\[
\pi_1+\frac12\pi_1+\frac76\pi_1=1.
\]

\[
\left(1+\frac12+\frac76\right)\pi_1=1.
\]

\[
\frac{16}{6}\pi_1=1
\quad\Longrightarrow\quad
\frac83\pi_1=1
\quad\Longrightarrow\quad
\pi_1=\frac38.
\]

Entonces:
\[
\pi_2=\frac12\cdot\frac38=\frac{3}{16},
\qquad
\pi_3=\frac76\cdot\frac38=\frac{7}{16}.
\]

Por lo tanto, la distribución estacionaria es
\[
\boxed{\pi=\left(\frac38,\frac{3}{16},\frac{7}{16}\right)}.
\]

\subsection*{¿También es distribución límite?}

Sí, porque la cadena es finita, irreducible y aperiódica. Entonces se cumple que la distribución estacionaria también es la distribución límite.

\[
\boxed{\text{Sí, también es distribución límite.}}
\]

\newpage
\section*{Ejercicio 7}

Considere la cadena con espacio de estados
\[
S=\{0,1,2,\dots\},
\]
y transiciones
\[
p_{i,i+1}=p,
\qquad
p_{i,0}=1-p,
\qquad i\geq 0,
\]
con
\[
\frac12<p<1.
\]

\subsection*{¿Es irreducible?}

Sí.

Desde cualquier estado \(i\), se puede ir a \(0\) en un paso con probabilidad \(1-p\). Además, desde \(0\) se puede llegar a cualquier estado \(j\) siguiendo
\[
0\to 1\to 2\to \cdots \to j,
\]
con probabilidad \(p^j>0\).

Entonces todos los estados se comunican y la cadena es irreducible.

\[
\boxed{\text{La cadena es irreducible.}}
\]

\subsection*{¿Es aperiódica?}

Sí, porque
\[
p_{00}=1-p>0.
\]

Entonces el estado \(0\) tiene periodo \(1\), y como la cadena es irreducible, todos los estados tienen periodo \(1\).

\[
\boxed{\text{La cadena es aperiódica.}}
\]

\subsection*{Distribución estacionaria}

Sea
\[
\pi=(\pi_0,\pi_1,\pi_2,\dots).
\]

Para que sea estacionaria, debe cumplirse que
\[
\pi=\pi P.
\]

Observando las ecuaciones de balance:
\[
\pi_1=p\pi_0,
\]
\[
\pi_2=p\pi_1,
\]
\[
\pi_3=p\pi_2,
\]
y en general
\[
\pi_j=p\pi_{j-1},\qquad j\geq 1.
\]

Entonces,
\[
\pi_j=p^j\pi_0.
\]

Ahora usamos que la suma de las probabilidades debe ser 1:
\[
\sum_{j=0}^{\infty}\pi_j=1.
\]

Sustituyendo:
\[
\sum_{j=0}^{\infty}p^j\pi_0=1.
\]

\[
\pi_0\sum_{j=0}^{\infty}p^j=1.
\]

Como la serie geométrica converge,
\[
\sum_{j=0}^{\infty}p^j=\frac{1}{1-p}.
\]

Entonces:
\[
\pi_0\frac{1}{1-p}=1
\quad\Longrightarrow\quad
\pi_0=1-p.
\]

Por lo tanto,
\[
\pi_j=(1-p)p^j,\qquad j\geq 0.
\]

Así, la distribución estacionaria es
\[
\boxed{\pi_j=(1-p)p^j,\qquad j=0,1,2,\dots}
\]

\subsection*{Observación}

Aunque el enunciado parece pedir que no exista distribución estacionaria, al hacer el desarrollo sí se obtiene una. Entonces, para la cadena tal como está escrita, sí existe distribución estacionaria.

\newpage
\section*{Ejercicio 8 [Punto extra]}

Con base en el ejercicio anterior, se pregunta si los estados son transitorios o nulo recurrentes cuando
\[
p>\frac12.
\]

Como en el ejercicio 7 encontramos una distribución estacionaria y la cadena es irreducible, entonces todos los estados son positivamente recurrentes.

Por lo tanto, no son ni transitorios ni nulo recurrentes.

\[
\boxed{\text{Todos los estados son positivamente recurrentes.}}
\]

\end{document}
